% Part: first-order-logic
% Chapter: natural-deduction
% Section: soundness

\documentclass[../../../include/open-logic-section]{subfiles}

\begin{document}

\iftag{FOL}
      {\olfileid{fol}{ntd}{sou}}
      {\olfileid{pl}{ntd}{sou}}
      
\olsection{السلامة}

\begin{explain}
نسق !!^a{derivation}، كالاستنتاج الطبيعي، \emph{سليم} متى امتنع عليه
!!{derive} ما لا يلزم في الحقيقة. فالسلامة ضمان لأنساق ال!!{derivation}
من مجاوزة ما تسوّغه المقدمات. ويتبيّن المراد بحسب الخاصية البرهانية
المبحوث عنها؛ فمن ذلك:
\begin{enumerate}
\item كل ما هو !!{derivable} من نوع !!{sentence} يكون \iftag{FOL}{صالحًا}{تحصيل حاصل}؛
\item إذا كانت !!a{sentence} !!{derivable} من جمل أخرى، فهي أيضًا
  نتيجة لها؛
\item إذا كانت مجموعة من ال!!{sentence}s غير متسقة، فهي غير قابلة
  للإرضاء.
\end{enumerate}
هذه خصائص لا غنى عنها لنسق !!a{derivation}؛ فإن تخلّفت واحدة منها
كان نسق ال!!{derivation} قاصرًا، إذ أمكنه !!{derive} ما لا حق له في
إثباته. فمن أجل ذلك كان إثبات سلامة نسق ال!!{derivation} من أوكد
ما يلزم في دراسته.
\end{explain}

\begin{thm}[السلامة]
\ollabel{thm:soundness}
إذا كانت~$!A$ !!{derivable} من الافتراضات ال!!{undischarged}
$\Gamma$، فإن~$\Gamma \Entails !A$.
\end{thm}

\begin{proof}
نأخذ~$\delta$، وهو !!a{derivation} لـ~$!A$، ونجعل عدد الاستدلالات
في~$\delta$ أساسًا للاستقراء.

أما الأساس، فعدد الاستدلالات فيه~$0$؛ فلا يشتمل~$\delta$ إلا على
!!{sentence} واحدة، هي~$!A$، مأخوذة افتراضًا. وهذا الافتراض
!!{undischarged}؛ فإن الافتراضات لا تصير !!{discharged} إلا باستدلال،
وليس هنا استدلال. فكل
\iftag{FOL}{!!{structure}~$\Struct{M}$}{!!{valuation}~$\pAssign{v}$}
ترضى فيها الافتراضات ال!!{undischarged} في البرهان ترضى فيها~$!A$
بالضرورة.

وأما الخطوة، فنفرض أن في~$\delta$ عددًا من الاستدلالات قدره~$n$.
وقد أمكن !!{derive} مقدمات الاستدلال الأخير بـ!!{derivation}s جزئية،
في كل منها أقل من~$n$ استدلالًا. فبحسب فرض
الاستقراء تلزم كل مقدمة من تلك المقدمات عن الافتراضات ال!!{undischarged}
في ال!!{derivation}s الجزئية التي انتهت إليها. ويبقى أن نبيّن لزوم
النتيجة~$!A$ عن الافتراضات ال!!{undischarged} في البرهان بأسره.

ولبيان ذلك نفصّل بحسب القاعدة التي ختم بها البرهان، ونبدأ بالقواعد
ذوات المقدمة الواحدة.
\begin{enumerate}
\item إن كان آخر الاستدلالات~\Intro{\lnot}، فصورة ال!!{derivation}
  هي:
  \begin{prooftree}
    \AxiomC{$\Gamma, \Discharge{!A}{n}$}
    \RightLabel{$\delta_1$}
    \DeduceC{$\lfalse$}
    \DischargeRule{\Intro{\lnot}}{n}
    \UnaryInfC{$\lnot !A$}
  \end{prooftree}
  بفرض الاستقراء يلزم~$\lfalse$ عن الافتراضات ال!!{undischarged}
  $\Gamma \cup \{!A\}$ في~$\delta_1$. ولنأخذ
  \iftag{FOL}{!!a{structure}~$\Struct{M}$}{!!a{valuation}~$\pAssign{v}$}.
  المطلوب أن يترتب على~$\iftag{FOL}{\Sat{M}}{\pSat{v}}{\Gamma}$
  صدق~$\iftag{FOL}{\Sat{M}}{\pSat{v}}{\lnot !A}$. فإن فرضنا للخلف
  $\iftag{FOL}{\Sat{M}}{\pSat{v}}{\Gamma}$ مع
  $\iftag{FOL}{\Sat/{M}}{\pSat/{v}}{\lnot !A}$، أي
  $\iftag{FOL}{\Sat{M}}{\pSat{v}}{!A}$، حصل
  $\iftag{FOL}{\Sat{M}}{\pSat{v}}{\Gamma \cup \{!A\}}$.
  ويلزم منه، بفرض الاستقراء، إرضاء الباطل، وهو محال. فلا بد من
  $\iftag{FOL}{\Sat{M}}{\pSat{v}}{\lnot !A}$.


\item وإن كان آخرها~\Elim{\land}، فإما أن نستنتج~$!A$
  وإما~$!B$ من المقدمة~$!A \land !B$. نبيّن الوجه الأول؛ وفيه
  ال!!{derivation}~$\delta$ على هذه الصورة:
  \begin{prooftree}
    \AxiomC{$\Gamma$}
    \RightLabel{$\delta_1$}
    \DeduceC{$!A \land !B$}
    \RightLabel{\Elim{\land}}
    \UnaryInfC{$!A$}
  \end{prooftree}
  بفرض الاستقراء تلزم~$!A \land !B$ عن الافتراضات
  ال!!{undischarged}~$\Gamma$ في~$\delta_1$. ولنأخذ
  \iftag{FOL}{!!a{structure}~$\Struct{M}$}{!!a{valuation}~$\pAssign{v}$}؛
  فينبغي أن يترتب على~$\iftag{FOL}{\Sat{M}}{\pSat{v}}{\Gamma}$
  صدق~$\iftag{FOL}{\Sat{M}}{\pSat{v}}{!A}$. فإذا فرضنا
  $\iftag{FOL}{\Sat{M}}{\pSat{v}}{\Gamma}$، كان فرض الاستقراء
  ($\Gamma \Entails !A \land !B$) مقتضيًا
  $\iftag{FOL}{\Sat{M}}{\pSat{v}}{!A \land !B}$. وتعريف الإرضاء
  يجعل~$\iftag{FOL}{\Sat{M}}{\pSat{v}}{!A \land !B}$ مكافئًا لاجتماع
  $\iftag{FOL}{\Sat{M}}{\pSat{v}}{!A}$ و
  $\iftag{FOL}{\Sat{M}}{\pSat{v}}{!B}$، وفي الأول المطلوب.
  وأما استنتاج~$!B$ من~$!A \land !B$ فبيانه على هذا النحو نفسه.


\item وإن كان آخرها~\Intro{\lor}، فاستنتاج
  $!A \lor !B$ يكون من~$!A$ أو من~$!B$. نأخذ الوجه الأول،
  وصورة ال!!{derivation} فيه:
  \begin{prooftree}
    \AxiomC{$\Gamma$}
    \RightLabel{$\delta_1$}
    \DeduceC{$!A$}
    \RightLabel{\Intro{\lor}}
    \UnaryInfC{$!A \lor !B$}
  \end{prooftree}
  بفرض الاستقراء تلزم~$!A$ عن الافتراضات ال!!{undischarged}
  $\Gamma$ في~$\delta_1$. ولنأخذ
  \iftag{FOL}{!!a{structure}~$\Struct{M}$}{!!a{valuation}~$\pAssign{v}$}.
  المطلوب أن يلزم من~$\iftag{FOL}{\Sat{M}}{\pSat{v}}{\Gamma}$
  صدق~$\iftag{FOL}{\Sat{M}}{\pSat{v}}{!A \lor !B}$. فإذا كان
  $\iftag{FOL}{\Sat{M}}{\pSat{v}}{\Gamma}$، كان
  $\iftag{FOL}{\Sat{M}}{\pSat{v}}{!A}$، إذ~$\Gamma \Entails !A$
  بفرض الاستقراء؛ ومن تعريف الفصل يلزم
  $\iftag{FOL}{\Sat{M}}{\pSat{v}}{!A \lor !B}$.
  وأما استنتاج~$!A \lor !B$ من~$!B$ فحجته نظير هذه.


\item وإن كان آخرها~\Intro{\lif}، فقد استنتجنا~$!A \lif !B$
  من برهان جزئي افتتح بافتراض~$!A$ وانتهى إلى~$!B$، هكذا:
  \begin{prooftree}
    \AxiomC{$\Gamma, \Discharge{!A}{n}$}
    \RightLabel{$\delta_1$}
    \DeduceC{$!B$}
    \DischargeRule{\Intro{\lif}}{n}
    \UnaryInfC{$!A \lif !B$}
  \end{prooftree}
  بفرض الاستقراء تلزم~$!B$ عن الافتراضات ال!!{undischarged} في
  $\delta_1$، أي~$\Gamma \cup \{!A\} \Entails !B$. ولنأخذ
  \iftag{FOL}{!!a{structure}~$\Struct{M}$}{!!a{valuation}~$\pAssign{v}$}.
  وليس من الافتراضات ال!!{undischarged} في~$\delta$ إلا~$\Gamma$،
  إذ سقط~$!A$ في الاستدلال الأخير. فالمطلوب
  $\Gamma \Entails !A \lif !B$. نفرض للخلف وجود
  \iftag{FOL}{!!{structure}~$\Struct{M}$}{!!{valuation}~$\pAssign{v}$}
  يتحقق فيه~$\iftag{FOL}{\Sat{M}}{\pSat{v}}{\Gamma}$، لكن
  $\iftag{FOL}{\Sat/{M}}{\pSat/{v}}{!A \lif !B}$؛
  أي إن~$\iftag{FOL}{\Sat{M}}{\pSat{v}}{!A}$ و
  $\iftag{FOL}{\Sat/{M}}{\pSat/{v}}{!B}$. غير أن~$!B$، بفرضنا،
  نتيجة لـ~$\Gamma \cup \{!A\}$؛ فيلزم
  $\iftag{FOL}{\Sat{M}}{\pSat{v}}{!B}$، وهو نقيض المفروض.
  فثبت~$\Gamma \Entails !A \lif !B$.


\item وإن كان آخرها~\FalseInt، فخاتمة~$\delta$ هي:
  \begin{prooftree}
    \AxiomC{$\Gamma$}
    \RightLabel{$\delta_1$}
    \DeduceC{$\lfalse$}
    \RightLabel{\FalseInt}
    \UnaryInfC{$!A$}
  \end{prooftree}
  بفرض الاستقراء~$\Gamma \Entails \lfalse$، والمطلوب
  $\Gamma \Entails !A$. فإن لم يكن، وجد
  $\iftag{FOL}{\Struct{M}}{\pAssign{v}}$ يتحقق فيه
  $\iftag{FOL}{\Sat{M}}{\pSat{v}}{\Gamma}$ و
  $\iftag{FOL}{\Sat/{M}}{\pSat/{v}}{!A}$. ولما كان
  $\iftag{FOL}{\Sat/{M}}{\pSat/{v}}{\lfalse}$ حاصلًا على كل حال،
  لزم~$\Gamma \Entails/ \lfalse$، مناقضًا فرض الاستقراء.


\item الاستدلال الأخير هو~\FalseCl: تمرين.

\iftag{FOL}{
\item وإن كان آخرها~\Intro{\lforall}، فصورة~$\delta$ هي:
  \begin{prooftree}
    \AxiomC{$\Gamma$}
    \RightLabel{$\delta_1$}
    \DeduceC{$!A(a)$}
    \RightLabel{\Intro{\lforall}}
    \UnaryInfC{$\lforall[x][!A(x)]$}
  \end{prooftree}
  المقدمة~$!A(a)$ تلزم عن الافتراضات ال!!{undischarged}~$\Gamma$
  بفرض الاستقراء. نأخذ بنية~$\Struct{M}$ تحقق~$\Sat{M}{\Gamma}$،
  ونبيّن~$\Sat{M}{\lforall[x][!A(x)]}$. ولما كانت
  $\lforall[x][!A(x)]$ !!a{sentence}، كفى أن نأخذ إسناد متغيرات
  كيفما اتفق~$s$ ونثبت~$\Sat{M}{!A(x)}[s]$
  (\olref[syn][ass]{prop:sat-quant}). وأعضاء~$\Gamma$ كلّها جمل؛
  فلدينا~$\Sat{M}{!B}[s]$ لكل~$!B \in \Gamma$، بحكم
  \olref[syn][sat]{defn:satisfaction}. نبني~$\Struct{M'}$ موافقةً
  لـ~$\Struct{M}$ في كل شيء سوى~$\Assign{a}{M'} = s(x)$.
  وحيث لا يرد~$a$ في~$\Gamma$، يبقى~$\Sat{M'}{\Gamma}$، بمقتضى
  \olref[syn][ext]{cor:extensionality-sent}. ومن
  $\Gamma \Entails !A(a)$ ينتج~$\Sat{M'}{!A(a)}$.
  ولأن~$!A(a)$ !!a{sentence}، فلدينا~$\Sat{M'}{!A(a)}[s]$ بحسب
  \olref[syn][ass]{prop:sentence-sat-true}. ثم إن
  $\Sat{M'}{!A(x)}[s]$ يكافئ~$\Sat{M'}{!A(a)}$، بمقتضى
  \olref[syn][ext]{prop:ext-formulas}؛ فإن~$!A(a)$ هي
  $\Subst{!A(x)}{a}{x}$. فقد ثبت~$\Sat{M'}{!A(x)}[s]$.
  ولعدم ورود~$a$ في~$!A(x)$، تنقلنا
  \olref[syn][ext]{prop:extensionality} إلى~$\Sat{M}{!A(x)}[s]$.
  ولم نخصّ~$s$ بشيء؛ فصحّ المطلوب لكل إسناد، ومن ثم
  $\Sat{M}{\lforall[x][!A(x)]}$.


\item الاستدلال الأخير هو~\Intro{\lexists}: تمرين.

\item الاستدلال الأخير هو~\Elim{\forall}: تمرين.
}{}
\end{enumerate}

وأما الاستدلالات ذوات المقدمات المتعددة فننظر فيها الآن، وهي:
\Elim{\lnot} و\Elim{\lor} و\Intro{\land} و\iftag{FOL}{\Elim{\lif} و
  \Elim{\lexists}}{\Elim{\lif}}.
\textbf{تنبيه تصحيحي على الأصل (C095-01).} أُعيدت $\Elim{\lnot}$ إلى قائمة
القواعد ذات المقدمات المتعددة؛ فهي مذكورة في تعداد الحالات اللاحق، وقد
سقط اسمها من هذه الإشارة التمهيدية وحدها.
\begin{enumerate}

\item إن كان آخرها~\Intro{\land}، فقد استنتجنا~$!A \land !B$
  من المقدمتين~$!A$ و$!B$، وصورة~$\delta$ هي:
  \begin{prooftree}
    \AxiomC{$\Gamma_1$}
    \RightLabel{$\delta_1$}
    \DeduceC{$!A$}
    \AxiomC{$\Gamma_2$}
    \RightLabel{$\delta_2$}
    \DeduceC{$!B$}
    \RightLabel{\Intro{\land}}
    \BinaryInfC{$!A \land !B$}
  \end{prooftree}
  بفرض الاستقراء تلزم~$!A$ عن الافتراضات ال!!{undischarged}
  $\Gamma_1$ في~$\delta_1$، كما تلزم~$!B$ عن الافتراضات
  ال!!{undischarged}~$\Gamma_2$ في~$\delta_2$.
  وأما الافتراضات ال!!{undischarged} في~$\delta$ فمجموعها
  $\Gamma_1 \cup \Gamma_2$؛ فيبقى إثبات
  $\Gamma_1 \cup \Gamma_2 \Entails !A \land !B$. نأخذ
  \iftag{FOL}{!!a{structure}~$\Struct{M}$}{!!a{valuation}~$\pAssign{v}$}
  تحقق~$\iftag{FOL}{\Sat{M}}{\pSat{v}}{\Gamma_1 \cup \Gamma_2}$.
  فمن~$\iftag{FOL}{\Sat{M}}{\pSat{v}}{\Gamma_1}$ يلزم
  $\iftag{FOL}{\Sat{M}}{\pSat{v}}{!A}$، إذ~$\Gamma_1 \Entails !A$؛
  ومن~$\iftag{FOL}{\Sat{M}}{\pSat{v}}{\Gamma_2}$ يلزم
  $\iftag{FOL}{\Sat{M}}{\pSat{v}}{!B}$، إذ~$\Gamma_2 \Entails !B$.
  وباجتماعهما نحصل على~$\iftag{FOL}{\Sat{M}}{\pSat{v}}{!A \land !B}$.


\item الاستدلال الأخير هو~\Elim{\lor}: تمرين.

\item وإن كان آخرها~\Elim{\lif}، فقد استنتجنا~$!B$
  من~$!A \lif !B$ و$!A$. فال!!{derivation}~$\delta$ صورته:
  \begin{prooftree}
    \AxiomC{$\Gamma_1$}
    \RightLabel{$\delta_1$}
    \DeduceC{$!A \lif !B$}
    \AxiomC{$\Gamma_2$}
    \RightLabel{$\delta_2$}
    \DeduceC{$!A$}
    \RightLabel{\Elim{\lif}}
    \BinaryInfC{$!B$}
  \end{prooftree}
  بفرض الاستقراء تلزم~$!A \lif !B$ عن الافتراضات ال!!{undischarged}
  $\Gamma_1$ في~$\delta_1$، وتلزم~$!A$ عن الافتراضات
  ال!!{undischarged}~$\Gamma_2$ في~$\delta_2$. نأخذ
  \iftag{FOL}{!!a{structure}~$\Struct{M}$}{!!a{valuation}~$\pAssign{v}$}.
  والمطلوب أن يترتب على
  $\iftag{FOL}{\Sat{M}}{\pSat{v}}{\Gamma_1 \cup \Gamma_2}$
  صدق~$\iftag{FOL}{\Sat{M}}{\pSat{v}}{!B}$. فإذا فرضنا
  $\iftag{FOL}{\Sat{M}}{\pSat{v}}{\Gamma_1 \cup \Gamma_2}$، لزم من
  $\Gamma_1 \Entails !A \lif !B$ صدق
  $\iftag{FOL}{\Sat{M}}{\pSat{v}}{!A \lif !B}$، ولزم من
  $\Gamma_2 \Entails !A$ صدق~$\iftag{FOL}{\Sat{M}}{\pSat{v}}{!A}$.
  فلا بد من~$\iftag{FOL}{\Sat{M}}{\pSat{v}}{!B}$؛ إذ لو كان
  $\iftag{FOL}{\Sat/{M}}{\pSat/{v}}{!B}$، لاجتمع ذلك مع
  $\iftag{FOL}{\Sat{M}}{\pSat{v}}{!A}$، فيلزم
  $\iftag{FOL}{\Sat/{M}}{\pSat/{v}}{!A \lif !B}$، مناقضًا
  $\iftag{FOL}{\Sat{M}}{\pSat{v}}{!A \lif !B}$.


\item الاستدلال الأخير هو~\Elim{\lnot}: تمرين.
  
\tagitem{FOL}{الاستدلال الأخير هو~\Elim{\lexists}: تمرين.}{}
\end{enumerate}
\end{proof}

\tagprob{FOL}
\begin{prob}
أكمل برهان \olref[fol][ntd][sou]{thm:soundness}.
\end{prob}
\tagendprob

\tagprob{notFOL}
\begin{prob}
أكمل برهان \olref[pl][ntd][sou]{thm:soundness}.
\end{prob}
\tagendprob

\begin{cor}
\ollabel{cor:weak-soundness}
إذا كان~$\Proves !A$، فإن~$!A$ \iftag{FOL}{صالحة}{تحصيل حاصل}.
\end{cor}

\begin{cor}
\ollabel{cor:consistency-soundness}
إذا كانت~$\Gamma$ قابلة للإرضاء، فهي متسقة.
\end{cor}

\begin{proof}
نأخذ المقابل بالنقيض. فإذا كانت~$\Gamma$ غير متسقة، كان
$\Gamma \Proves \lfalse$؛ أي وجد !!a{derivation} لـ~$\lfalse$،
افتراضاته ال!!{undischarged} في~$\Gamma$. وبمبرهنة
\olref{thm:soundness} يلزم أن تكون كل
\iftag{FOL}{!!{structure}~$\Struct{M}$}{!!{valuation}~$\pAssign{v}$}
ترضي~$\Gamma$ مرضيةً لـ~$\lfalse$ أيضًا. غير أن
$\iftag{FOL}{\Sat/{M}}{\pSat/{v}}{\lfalse}$ في كل
\iftag{FOL}{!!{structure}~$\Struct{M}$}{!!{valuation}~$\pAssign{v}$}.
فلا توجد \iftag{FOL}{$\Struct{M}$}{$\pAssign{v}$} ترضي~$\Gamma$؛
ومعناه أن~$\Gamma$ غير قابلة للإرضاء.
\end{proof}

\end{document}
